\section{Problems}

These follow the same order as the chapter. Hints for the starred ones are at the end.

\begin{prob}
Write the permutation $(1\,3)(2\,5\,4)$ as a product of neighbour swaps, and check your answer by
composing them one at a time. Then do the same for an arbitrary $k$-cycle, and deduce that a
$k$-cycle is a product of $k-1$ swaps.
\end{prob}

\begin{prob}
The regular tetrahedron has four corners. Show that its rotation group has $12$ elements, and that
allowing reflections gives a group isomorphic to $S_4$. Which groups arise this way from the other
regular solids? Do the higher-dimensional cubes give symmetric groups?$^\star$
\end{prob}

\begin{prob}
Show that $m$ generates $\Z/n\Z$ if and only if $m$ and $n$ have no common factor greater than
$1$.$^\star$ Then list all subgroups of $\Z/12\Z$ and match them to the divisors of $12$.
\end{prob}

\begin{prob}
Construct a group of matrices all of whose entries lie in a finite set.$^\star$ How many elements
does your group have?
\end{prob}

\begin{prob}
Show that every word in $a,b,a\inv,b\inv$ reduces to exactly one reduced word.$^\star$ Where in the
construction of $F_2$ is this needed?
\end{prob}

\begin{prob}
Let $u$ and $v$ be reduced words. Show that the reduction of the concatenation $uv$ deletes a block
of letters at the seam and nothing else, and use this to compute the length of $uv$ in terms of the
lengths of $u$ and $v$ and the size of the cancelling block.
\end{prob}

\begin{prob}
Prove the cancellation laws: in any group, $gx=gy$ implies $x=y$, and $xg=yg$ implies $x=y$. Give an
example of a set with a multiplication satisfying (G1) and (G3) in which cancellation fails.
\end{prob}

\begin{prob}
Show that a group in which every element satisfies $g^2=1$ is abelian. Then find such a group with
exactly $8$ elements.
\end{prob}

\begin{prob}
Let $f:G\to H$ be a homomorphism. Show that the image of $f$ is a subgroup of $H$ and the kernel is
a normal subgroup of $G$. Show that $f$ is injective exactly when its kernel is $\{1\}$.
\end{prob}

\begin{prob}\label{prob:cayleyD3}
Draw the Cayley graph of $D_3$ with respect to $\{s,t\}$, its two reflections, and with respect to
$\{r,s\}$, a rotation and a reflection. Compare the two pictures.
\end{prob}

\begin{prob}
Show that $\langle a,b\mid aba\inv b\inv\rangle\cong\Z^2$ by proving that every word has a unique
normal form $a^mb^n$.$^\star$
\end{prob}

\begin{prob}
Show that $\langle a\mid a^n\rangle\cong\Z/n\Z$. Then decide what group
$\langle a,b\mid a^2, b^3, (ab)^3\rangle$ is, by drawing enough of its Cayley graph to recognise it.
\end{prob}

\subsection*{Hints}

\noindent\textbf{2.} Count as in Proposition~\ref{prop:24}: a corner can go to any of four corners,
and then three rotations fix that choice. For the last part, count the rotations of the
four-dimensional cube and compare with $5!$ and $6!$; the numbers do not match.

\noindent\textbf{3.} If $\gcd(m,n)=1$ there are integers $u,v$ with $um+vn=1$; reduce modulo $n$.
For the converse, note that every multiple of $m$ is a multiple of any common divisor.

\noindent\textbf{4.} Use entries in $\Z/p\Z$ for a prime $p$, where every nonzero element has a
multiplicative inverse, and take the invertible $2\times2$ matrices. For $p=2$ there are six.

\noindent\textbf{5.} Induct on length, as in Theorem~\ref{thm:unique}: two available deletions
either do not overlap or overlap in one letter.

\noindent\textbf{11.} The relator lets you swap a $b$ past an $a$. Show the swapping terminates, and
that the resulting $a^mb^n$ cannot be changed further.
